\paragraph{Second-order response in the exterior.} The \(f6\) row requires \(m^2q_i\), since the
first-order response in \(y_i/p_i\) cancels with \(t_{ic}\). To verify the second-order source, set
\(X=(a_{i,I}^{\mathsf T})_{i\in E'}\), \(D_p=\operatorname{diag}(p_i)_{i\in E'}\). Subtracting the
stationarity equations gives \(y_i/p_i=(X\eta)_i+(B_1)_i\), where \(\eta_f=\gamma_f/h_f\) on the
diagonals and the other components are row and column multipliers. With \(\theta_i=y_i/p_i\), subtracting
the logarithms of the stationarity equations gives the exact identity \[
 (B_1)_i=\theta_i-\log(1+\theta_i)-3\log(1-y_i/(1-p_i))
 +\sum_{f\in \familias(i)\cap(D\cup A)}[\log(1+\eta_f)-\eta_f].
\] The mass equations give \(\mathcal M\eta=d_0-X^{\mathsf T}D_pB_1\). Hence \[
 q=B_1+X\mathcal M^+(r-X^{\mathsf T}D_p B_1).
\] Taylor expansion in the squares and vacancies gives \(m^2B_1=L+O_b(m^{-1})\), with \[
 L_i=\tfrac12(mt_{ic})^2-\frac{3mp_i(mt_{ic})}{1-p_i}
       -\tfrac12\sum_{f\in \familias(i)\cap(D\cup A)}(mz_f^{\rm phys})^2,\qquad
 z^{\rm phys}=(Ra_c)_I .
\] Here we used \(y_i/p_i=-t_{ic}+O_b(m^{-2})\) and \(\gamma_f/h_f=-z_f^{\rm phys}+O_b(m^{-2})\). On the
complement of the null mode, the block form of \(S_0R=I-\Pi_m\) gives \(\|\mathcal
M^+-R_{II}\|_\infty=O_b(m^{-1})\) in the maximum absolute row-sum norm, since the retained lines lose
mass \(O_b(m^{-1})\) and the inverse is bounded. Thus \(m^2q=L+XR_{II}(m^2r-X^{\mathsf
T}D_pL)+O_b(m^{-1})\).

We express the contraction using the same profile kernels. Fix the center \(c\) and put
\(\Lambda(i)=T_i^2/2-3W(i)T_i-\tfrac12\sum_{f\in \familias(i)\cap(D\cup A)}z_f^2\). For a retained line
\(u\), write \begin{align*}
 \omega(u)={}&-\sum_{j\in \familias(c)} \pesodirecto(u\cap j)
     \left[t_j+\sum_{a\in \familias(u\cap j)\setminus\{j\}}z_a\right]
      -\int_{i'\in u}\mathbb E\Lambda(i')\,d\omega_u(i'),\\
 q(i)={}&\Lambda(i)+\sum_{u\in \familias(i)}
 \left[\omega(u)+\int E(u,u')\omega(u')\,d\rho(u')\right],\\
 f6={}&-\left\langle w(i)\mathbb E q(i)\right\rangle_{i,c}.
\end{align*} The first sum omits missing intersections. The integrated fiber term in \(\omega\) is the
mean with \(i'\) marked and \(c\) fixed. In the products in the sum over intersections, \(\psi\) applies
to the whole monomial. The two parts of \(\omega\) come from \(m^2r=-m^2(C_0t+U_0z)\) and \(-X^{\mathsf
T}D_pL\). The identity and integral parts of \(R_{II}\) give the two terms in brackets in \(q(i)\). The
area mean carries the sign of \(-m^2\sum_{E'}p_iq_iw_i\). Lemma~\ref{lem:f6} proves the trace
rearrangement, including the walk in Figure~\ref{fig:rectangulos}.
